\documentclass[10pt,spanish]{article}
\usepackage[T1]{fontenc}
\usepackage[utf8]{inputenc}
\usepackage[a4paper]{geometry}
\usepackage{verbatim}
\geometry{verbose,tmargin=1.5cm,bmargin=0.5cm,lmargin=0.7cm,rmargin=0.7cm}
\usepackage{amsmath,amssymb,color,fancyhdr,mathtools,multicol}
\setlength{\columnsep}{0.7cm}
%\usepackage{fontspec}
%\setmainfont{Arial}

\makeatletter
\newcommand{\mathcircumflex}[0]{\mbox{\^{}}}
\newcommand{\myunder}[1]{\underline{\textbf{#1}}}

\usepackage{xcolor}
\definecolor{codegreen}{rgb}{0,0.6,0}                                                          
\definecolor{codegray}{rgb}{0.5,0.5,0.5}                                                       
\definecolor{codepurple}{rgb}{0.58,0,0.82}                                                     
\definecolor{backcolour}{rgb}{0.95,0.95,0.92}
\makeatother
\usepackage{babel}
\addto\shorthandsspanish{\spanishdeactivate{~<>}}

\usepackage{listings}
\lstset{language={C++},
basicstyle={\ttfamily},
breaklines=true,
tabsize=2,
backgroundcolor={\color{backcolour}},
commentstyle={\color{codegreen}},
keywordstyle={\color{magenta}},
numberstyle={\color{codegray}},
stringstyle={\color{codepurple}},
%breakatwhitespace=false,
captionpos=b,
keepspaces=true,
numbers=left,
numbersep=5pt,%
}

\addto\captionsenglish{\renewcommand{\lstlistingname}{Listing}}
\addto\captionsspanish{\renewcommand{\lstlistingname}{Listado de código}}
\renewcommand{\lstlistingname}{Listado de código}

\begin{document}
%\small
\pagestyle{fancy}
\thispagestyle{fancy}
\fancyhead[L]{Universidad de Murcia---Magos del tle}
\fancyhead[R]{Page \thepage}
\begin{multicols}{2}

\tableofcontents

\section{Number Theory}

\subsection{Criba de Gries-Misra $\mathcal{O}(N)$}

\begin{lstlisting}
vector<int> mindiv;
vector<int> primes;

void sieve(int N) {
    mindiv = vector<int>(N+1, 0);
    primes = vector<int>();
    for(int i=2; i <= N; ++i) {
        if(mindiv[i] == 0) {
            mindiv[i] = i;
            primes.push_back(i);
        }
        for(int j=0; j < primes.size() and primes[j] <= mindiv[i] and primes[j]*i <= N; ++j) {
            mindiv[i*primes[j]] = primes[j];
        }
    }
}
\end{lstlisting}

\subsection{Discrete Log $\mathcal{O}(\sqrt{m})$}

Se basa en verificar si la ecuación en congruencias $a^x \equiv b \mod m$ tiene solución y, en caso afirmativo, dar el valor de tal $x$. La idea clave es expresar $x = nq - p$, donde $n$ nos interesará tomarlo lo más próximo posible a $\sqrt{m}$.

\section{Estructuras de Datos}

\subsection{Fenwick Tree $\mathcal{O}(\log m)$}

\begin{lstlisting}
void update(int i) {
    while (i <= m) ++st[i], i += i&-i;
}

int query(int i) {
    int s = 0;
 
    while (i >= 1) s += st[i], i -= i&-i;
 
    return s;
}
\end{lstlisting}

\section{Grafos}

\subsection{Eulerian Path/Circuit}

\begin{lstlisting}
stack<int> st; st.push(1);
vector<int> pth;
while (st.size()) {
    int u = st.top();
    bool f = false;
    while (G[u].size()) {
        int v = G[u].back().first, i = G[u].back().second; G[u].pop_back();
        if (!vis[i]) {
            st.push(v);
            vis[i] = true, f = true;
            break;
        }
    }
    if (!f) pth.push_back(u), st.pop();
}
\end{lstlisting}

\subsection{Max-Flow $\mathcal{O}(m^2\log c)$}

\begin{lstlisting}
bool dfs(int u, int mini, int k) {
    if (vis[u]) return false;
    if (u == n) {
        val = mini;
        return true;
    }
 
    vis[u] = true;
 
    for (int v = 1; v <= n; ++v) if (G[u][v] >= k) {
        if (dfs(v, min(mini, G[u][v]), k)) {
            G[u][v] -= val;
            G[v][u] += val;
            return true;
        }
    }
 
    return false;
}

int ans = 0;
for (int k = 1e9; k; k >>= 1) {
    vis = vector<bool>(n + 1, false);
    while (dfs(1, 1e9 + 5, k)) {
        ans += val;
        vis = vector<bool>(n + 1, false);
    }
}
\end{lstlisting}

\section{Strings}

\subsection{Z-algorithm $\mathcal{O}(n)$}

\begin{lstlisting}
int n = s.size();
vector<int> z(n);
int x = 0, y = 0;

for (int i = 1; i < n; i++) {
    z[i] = max(0, min(z[i - x], y - i + 1));

    while (i + z[i] < n && s[z[i]] == s[i + z[i]]) {
        x = i, y = i + z[i], z[i]++;
    }

    if (i > m && z[i] == m) ans++; 
}
\end{lstlisting}

\subsection{Suffix Array}

\begin{lstlisting}
int n = s.size();
 
s += '$';

vector<iii> a;
for (int i = 0; i <= n; ++i) a.push_back({{s[i], 0}, i});

sort(a.begin(), a.end());

for (int k = 1; k <= n; k <<= 1) {
    vector<int> p(n + 1);

    p[a[0].second] = 0;

    for (int i = 1; i <= n; ++i) {
        p[a[i].second] = p[a[i - 1].second];
        if (a[i].first != a[i - 1].first) ++p[a[i].second];
    }

    for (int i = 0; i <= n; ++i) a[i] = {{p[i], p[(i + k) % (n + 1)]}, i};

    sort(a.begin(), a.end());
}

vector<int> p(n + 1);
for (int i = 0; i <= n; ++i) p[a[i].second] = i;

int k = 0;
vector<int> lcp(n + 1, 0);
for (int i = 0; i < n; ++i) {
    int j = a[p[i] - 1].second;

    k = max(0, k - 1);

    while (i + k < n && j + k < n && s[i + k] == s[j + k]) ++k;

    lcp[p[i]] = k;
}
\end{lstlisting}

\end{multicols}
%\pagebreak~~
\end{document}